\section{Conclusion}


In this paper, we have proposed a unified dual-mask physical model to address the persistent challenge of depth estimation for non-Lambertian surfaces in light field imaging, where the fundamental Epipolar Plane Image (EPI) linearity assumption is severely violated. By explicitly modeling the physical degradation of angular signals, our framework decouples material-specific reflections from geometric structures, enabling robust depth inference in complex mixed-reflectance environments. 

The main contributions of this work are summarized as follows:
\textit{ \textbf{Geometry-based Material Classification via Angular Signal Decomposition:} We introduced a three-layer decomposition mechanism that overcomes frequency aliasing caused by discrete low-resolution angular sampling, providing a robust and efficient material identification scheme that lays the foundation for dual-branch depth estimation.
} \textbf{Unified Dual-Mask Depth Estimation Architecture:} We developed a dual-branch framework that effectively mitigates the depth collapse typically encountered by single-EPI-based methods when facing specular or translucent surfaces due to physical assumption failures.
\textbf{Mask Routing for Mixed-Scene Generalization:} By integrating a dynamic mask routing mechanism within a unified architecture, the proposed model achieves high-precision depth estimation across both mixed and holistic scenes, significantly improving the generalization capability and practical applicability of the system.

Although the proposed unified architecture achieves robust performance in mixed scenarios, it remains fundamentally constrained by the extreme scarcity of non-Lambertian training data and the inherent blurring of EPI slopes in complex textured regions. To transcend these intrinsic bottlenecks of epipolar geometry, future research will shift towards non-EPI paradigms, such as Neural Radiance Fields (NeRF), leveraging substantially expanded multi-reflectance datasets to achieve comprehensive depth estimation across all material domains.


The explicit three-layer angular signal decomposition offers a distinct advantage over purely data-driven, black-box convolutional neural networks (CNNs) by enforcing physical priors directly within the signal processing pipeline. By decoupling the angular signal $I(u,v)$ into the DC component (ambient illumination), the linear component (epipolar geometry), and the residual $\epsilon(u,v)$ (material reflectance), our model prevents the entanglement of high-frequency specular highlights with low-frequency disparity cues. In end-to-end architectures, the network often struggles to disentangle these overlapping frequency bands, leading to depth collapse in non-Lambertian regions where the Bidirectional Reflectance Distribution Function (BRDF) varies rapidly. Our physics-driven decoupling ensures that the residual $\epsilon(u,v)$ safely absorbs material-induced radiometric variations, guaranteeing that the geometric disparity estimation remains strictly invariant to surface reflectance properties, thereby enhancing both algorithmic robustness and physical interpretability.

A critical design choice in our framework is the deliberate departure from frequency-domain analyses, such as the 2D Discrete Fourier Transform (DFT), in favor of spatial-domain geometric features. From a signal sampling perspective, commercial plenoptic cameras typically provide low angular resolutions (e.g., $9 \times 9$ views). According to the Nyquist-Shannon sampling theorem, such sparse angular sampling inevitably induces severe spectral aliasing and leakage in the frequency domain, rendering DFT-based material classification fundamentally unreliable. By extracting spatial geometric features—such as angular gradients, correlation coefficients, and coefficients of variation—we effectively bypass the frequency aliasing bottleneck. These geometric descriptors act as robust spatial-domain statistics that preserve structural integrity without requiring dense angular sampling. The high Cohen's $d$ effect size of the angular gradient feature empirically validates that spatial-domain metrics remain highly discriminative for material identification even under severely discrete angular sampling conditions.

Furthermore, the proposed unified dual-mask routing mechanism addresses the inherent gradient conflict encountered when optimizing a single monolithic network for mixed-reflectance scenes. In conventional single-branch models, the loss gradients derived from Lambertian and non-Lambertian regions often point in opposing directions in the parameter space, causing sub-optimal convergence and feature suppression. By dynamically routing features to specialized sub-networks based on the material masks, we effectively isolate the optimization landscapes. Computationally, utilizing a lightweight Random Forest (RF) classifier for the initial mask generation is significantly more efficient than deploying heavy, pixel-wise semantic segmentation networks. However, this cascaded architecture introduces a potential vulnerability: misclassifications in the RF stage could propagate as erroneous hard masks. To mitigate this, the domain-balanced sampling strategy and the inherent tolerance of the multi-directional EPI features ensure that minor mask boundary inaccuracies do not catastrophically degrade the final depth maps. Future iterations could explore differentiable soft-mask routing to enable unified end-to-end gradient flow.

Despite the demonstrated efficacy, the current physical model operates under specific boundary conditions. The linear approximation of the epipolar component assumes relatively small baselines and opaque or semi-translucent surfaces. For highly refractive transparent objects (e.g., thick glass or fluids), light rays undergo complex bending, fundamentally violating the straight-line EPI assumption and producing highly non-linear, fragmented trajectories. Addressing such extreme optical phenomena would require integrating explicit ray-tracing priors or transitioning to volumetric rendering paradigms, such as Neural Radiance Fields (NeRF) or 3D Gaussian Splatting, specifically adapted for multi-view light field inputs. Additionally, extending the current discrete angular decomposition to continuous light field representations could unlock higher-fidelity depth estimation for next-generation plenoptic sensors, while incorporating temporal consistency constraints would be essential for deploying this framework in dynamic, real-world robotic navigation scenarios.

While the proposed unified dual-mask framework demonstrates robust performance in decoupling angular signals, several physical and computational limitations remain, outlining critical directions for future research.

First, the three-layer angular signal decomposition relies on a first-order local linearity assumption, modeled as $I(u,v) = I_{dc} + a \cdot u + b \cdot v + \epsilon(u,v)$. While this effectively isolates material residuals $\epsilon(u,v)$ for standard specular and scattering surfaces, it exhibits limitations under complex global illumination, such as inter-reflections or caustics. In such regimes, the residual term becomes non-linearly coupled with geometric disparity, potentially degrading the Random Forest classifier's precision and inducing mask routing artifacts at high-contrast material boundaries. Future work could integrate higher-order bidirectional reflectance distribution function (BRDF) models or hybridize our explicit physical decomposition with implicit neural rendering to better disentangle multi-bounce light transport.

Second, although our geometry-based feature extraction circumvents frequency aliasing at low angular resolutions (e.g., $9 \times 9$), the discrete spatial sampling of the Epipolar Plane Image (EPI) restricts the recovery of sub-pixel or extremely fine structures (e.g., hair, thin wires). At these scales, EPI linearity inherently degrades due to spatial aliasing, leading to potential depth discontinuities. To mitigate this, future iterations could incorporate spatial super-resolution priors or exploit sub-pixel phase correlation to enhance EPI structural integrity prior to dual-mask routing.

Finally, the current signal processing pipeline is strictly formulated for static light fields. Extending the architecture to spatiotemporal light fields introduces temporal motion blur, which acts as an additional low-pass filter on angular signals, further obscuring material-specific high-frequency residuals. Future research will investigate integrating temporal consistency constraints into the mask routing mechanism. Furthermore, from a hardware-algorithm co-design perspective, leveraging polarized light field sensors or event-based vision could provide direct physical measurements of surface normals and material properties. This would bypass the computational overhead of algorithmic signal decoupling, enabling real-time, high-fidelity depth estimation in dynamic, mixed-reflectance environments.